% GENERATED FILE. Edit canonical lessons, then run npm run generate:books.
% canonical-source-hash: fnv1a64:9cbf0327c85b604c
% canonical-lessons: GE-C73-traube, GE-C73-marmelade, GE-C73-kaese, GE-C73-salz, GE-C73-zucker

\chapter{Grape, Jam, Cheese, Salt, Sugar}
\label{ch:ge-grape-jam-cheese-salt-sugar}
\hlchaptermodality{73}

\begin{chapteropening}
\textbf{By the end of this chapter:} \emph{I can say a grape, jam, cheese, salt and sugar.}

Complete the last lesson of chapter 73: I can say a grape, jam, cheese, salt and sugar.
\end{chapteropening}

\section[die Traube]{die Traube — a grape}

\label{lesson:GE-C73-traube}



\begin{quote}
Before the new one: say the German for a pear, then the German for a lemon.
\end{quote}

\subsection*{You'll want to know: die Traube}
\textbf{die Traube} — \textquotedblleft{}a grape\textquotedblright{}.

\begin{cousinweb}
\textbf{Traube} $\leftarrow$ Old High German \emph{drūbo} $\leftarrow$ Proto-Germanic \emph{*þrūbô}, \textquotedblleft{}grape\textquotedblright{} — inherited; Dutch \emph{druif} is the same word. (EtymDB 2.1)
\end{cousinweb}

\begin{grammarlens}[title={What you've built}]
One more word for this chapter.
\end{grammarlens}

\subsection*{Your turn}
\begin{itemize}
\raggedright
  \item \emph{Say it:} \emph{die Traube}
  \item \emph{Say it:} \emph{die Traube}, once more
  \item \emph{Say it:} say \emph{die Traube}
  \item \emph{Recall:} say the German for a pear, then the German for a lemon, then say \emph{die Traube} again
\end{itemize}

\begin{tcolorbox}[breakable,colback=teal!4,colframe=teal!35!black,title={Before you move on}]
What does die Traube mean? (\textquotedblleft{}A grape\textquotedblright{}.) Say it once more.
\end{tcolorbox}
\section[die Marmelade]{die Marmelade — jam}

\label{lesson:GE-C73-marmelade}



\begin{quote}
Before the new one: say the German for a lemon, then the German for a grape.
\end{quote}

\subsection*{You'll want to know: die Marmelade}
\textbf{die Marmelade} — \textquotedblleft{}jam\textquotedblright{}.

\begin{cousinweb}
\textbf{Marmelade} $\leftarrow$ French \emph{marmelade} $\leftarrow$ Portuguese \emph{marmelada}, \textquotedblleft{}quince jam\textquotedblright{}, from \emph{marmelo}, \textquotedblleft{}quince\textquotedblright{} (EtymDB 2.1) — borrowed.
\end{cousinweb}

\begin{grammarlens}[title={What you've built}]
One more word for this chapter.
\end{grammarlens}

\subsection*{Your turn}
\begin{itemize}
\raggedright
  \item \emph{Say it:} \emph{die Marmelade}
  \item \emph{Say it:} \emph{die Marmelade}, once more
  \item \emph{Say it:} say \emph{die Marmelade}
  \item \emph{Recall:} say the German for a lemon, then the German for a grape, then say \emph{die Marmelade} again
\end{itemize}

\begin{tcolorbox}[breakable,colback=teal!4,colframe=teal!35!black,title={Before you move on}]
What does die Marmelade mean? (\textquotedblleft{}Jam\textquotedblright{}.) Say it once more.
\end{tcolorbox}
\section[der Käse]{der Käse — cheese}

\label{lesson:GE-C73-kaese}



\begin{quote}
Before the new one: say the German for a grape, then the German for jam.
\end{quote}

\subsection*{You'll want to know: der Käse}
\textbf{der Käse} — \textquotedblleft{}cheese\textquotedblright{}.

\begin{cousinweb}
\textbf{Käse} $\leftarrow$ Old High German \emph{chāsi} $\leftarrow$ Latin \emph{cāseus}, \textquotedblleft{}cheese\textquotedblright{} (EtymDB 2.1) — an early loan, as English \emph{cheese} was.
\end{cousinweb}

\begin{grammarlens}[title={What you've built}]
One more word for this chapter.
\end{grammarlens}

\subsection*{Your turn}
\begin{itemize}
\raggedright
  \item \emph{Say it:} \emph{der Käse}
  \item \emph{Say it:} \emph{der Käse}, once more
  \item \emph{Say it:} say \emph{der Käse}
  \item \emph{Recall:} say the German for a grape, then the German for jam, then say \emph{der Käse} again
\end{itemize}

\begin{tcolorbox}[breakable,colback=teal!4,colframe=teal!35!black,title={Before you move on}]
What does der Käse mean? (\textquotedblleft{}Cheese\textquotedblright{}.) Say it once more.
\end{tcolorbox}
\section[das Salz]{das Salz — salt}

\label{lesson:GE-C73-salz}



\begin{quote}
Before the new one: say the German for jam, then the German for cheese.
\end{quote}

\subsection*{You'll want to know: das Salz}
\textbf{das Salz} — \textquotedblleft{}salt\textquotedblright{}.

\begin{cousinweb}
\textbf{Salz} $\leftarrow$ Old High German \emph{salz} $\leftarrow$ Proto-Germanic \emph{*saltą} — inherited, the same word as English \emph{salt}, with German \emph{z} where English kept \emph{t}. (EtymDB 2.1)
\end{cousinweb}

\begin{grammarlens}[title={What you've built}]
One more word for this chapter.
\end{grammarlens}

\subsection*{Your turn}
\begin{itemize}
\raggedright
  \item \emph{Say it:} \emph{das Salz}
  \item \emph{Say it:} \emph{das Salz}, once more
  \item \emph{Say it:} say \emph{das Salz}
  \item \emph{Recall:} say the German for jam, then the German for cheese, then say \emph{das Salz} again
\end{itemize}

\begin{tcolorbox}[breakable,colback=teal!4,colframe=teal!35!black,title={Before you move on}]
What does das Salz mean? (\textquotedblleft{}Salt\textquotedblright{}.) Say it once more.
\end{tcolorbox}
\section[der Zucker]{der Zucker — sugar}

\label{lesson:GE-C73-zucker}



\begin{quote}
Before the new one: say the German for cheese, then the German for salt.
\end{quote}

\subsection*{You'll want to know: der Zucker}
\textbf{der Zucker} — \textquotedblleft{}sugar\textquotedblright{}.

\begin{cousinweb}
\textbf{Zucker} $\leftarrow$ Italian \emph{zucchero} $\leftarrow$ Arabic \emph{sukkar} (EtymDB 2.1) — borrowed with the trade in sugar.
\end{cousinweb}

\begin{grammarlens}[title={What you've built}]
One more word for this chapter.
\end{grammarlens}

\subsection*{Your turn}
\begin{itemize}
\raggedright
  \item \emph{Say it:} \emph{der Zucker}
  \item \emph{Say it:} \emph{der Zucker}, once more
  \item \emph{Say it:} say \emph{der Zucker}
  \item \emph{Recall:} say the German for cheese, then the German for salt, then say \emph{der Zucker} again
\end{itemize}

\begin{tcolorbox}[breakable,colback=teal!4,colframe=teal!35!black,title={Before you move on}]
What does der Zucker mean? (\textquotedblleft{}Sugar\textquotedblright{}.) Say it once more.
\end{tcolorbox}
